% ==============================================================================
% CHƯƠNG 2: CƠ SỞ LÝ THUYẾT
% ==============================================================================
\chapter{Cơ sở lý thuyết}
\label{chap:co-so}

\section{Các khái niệm nền tảng}
\label{sec:khai-niem}

\begin{dinhnghia}[Mô hình dựa trên tác vụ]
  \label{dn:tac-vu}
  Một công việc được xem là tập hợp các \textit{tác vụ} riêng biệt; công nghệ có thể
  thay thế một phần tác vụ mà không thay thế toàn bộ công việc.
\end{dinhnghia}

Theo Định nghĩa~\ref{dn:tac-vu}, tác động của công nghệ được tách thành hai hiệu ứng:
\begin{itemize}
  \item \textbf{Hiệu ứng thay thế:} công nghệ làm tác vụ rẻ hơn, giảm nhu cầu lao động
        ở tác vụ đó.
  \item \textbf{Hiệu ứng năng suất:} chi phí giảm giúp mở rộng sản xuất, tăng nhu cầu
        lao động ở các tác vụ còn lại.
\end{itemize}

\section{Mô hình toán học}
\label{sec:mo-hinh}

Tác động ròng lên nhu cầu lao động có thể viết gọn như \PT{eq:tac-dong}:
\begin{khungcongthuc}[Tác động ròng]
  \begin{equation}
    \label{eq:tac-dong}
    \Delta L \;=\; \underbrace{-\,\alpha \, \Delta A}_{\text{thay thế}}
          \;+\; \underbrace{\beta \, \Delta A}_{\text{năng suất}},
    \qquad \alpha, \beta > 0
  \end{equation}
\end{khungcongthuc}
trong đó $\Delta A$ là mức thay đổi tiếp xúc với công nghệ. Kết quả cần nhớ có thể
đóng khung ngay trong dòng công thức:
\begin{equation}
  \label{eq:dieu-kien}
  \NoiBat{\Delta L < 0 \iff \alpha > \beta}
\end{equation}

\begin{dinhly}[Điều kiện cân bằng]
  Nếu $\alpha = \beta$ thì tác động ròng bằng $0$ dù từng hiệu ứng thành phần khác $0$.
\end{dinhly}

\begin{vidu}[Áp dụng công thức]
  Cho $\alpha = 0{,}6$, $\beta = 0{,}4$ và $\Delta A = 10$. Tính $\Delta L$.
\end{vidu}
\begin{loigiai}
  Thay số vào \PT{eq:tac-dong}:
  $\Delta L = -0{,}6 \cdot 10 + 0{,}4 \cdot 10 = -2$.
  Vì $\alpha > \beta$ nên nhu cầu lao động giảm, đúng với điều kiện \eqref{eq:dieu-kien}.
\end{loigiai}

\section{Các nghiên cứu liên quan}
\label{sec:lien-quan}

Tóm lược các công trình chính \cite{acemoglu2022, brynjolfsson2025, acemoglu2020}, chỉ ra
điểm giống/khác với hướng tiếp cận của báo cáo.

\begin{trichdan}{Acemoglu và cộng sự (2022)}
  AI technologies are still in their infancy and have spread to only a limited part of
  the US economy.
\end{trichdan}

\begin{ychinh}
  \begin{dssao}
    \item Ý chính thứ nhất của chương.
    \item Ý chính thứ hai, dẫn dắt sang chương sau.
  \end{dssao}
\end{ychinh}
